\section{Immutable Source Snapshot Layer}

Before a source record is transformed into a metric observation, the system
preserves the retrieved response as an immutable source snapshot. This is a
vertical-neutral storage contract: the same structure can contain an SEC
filing for an NCRA, an inspection report for a corrections department, or a
court decision for a judiciary profile.

\begin{lstlisting}
entity_id + source_type + source_record_id
  -> retrieved payload
  -> canonical content hash
  -> immutable snapshot
  -> parser/versioned observation
\end{lstlisting}

Each snapshot records the entity, source type, source identifier, retrieval
time, optional source URL, parser version, payload, and content hash. The
content hash permits integrity verification and prevents a later source
change from silently altering a previously computed result. A parser may
produce a new observation version, but it must not mutate the original
snapshot.

Snapshot integrity is checked before selection for downstream processing.
Invalid or tampered snapshots are excluded rather than silently repaired.
This establishes the reproducibility boundary required by the Beacon baseline:
auditers can re-run a parser against the preserved source response and compare
both the source hash and the resulting observation.

The current implementation provides deterministic in-memory contracts. Durable
object storage, retention policy, encrypted storage for restricted material,
source-specific fetch adapters, and signed external manifests remain later
infrastructure work. No snapshot by itself constitutes a validated metric or
published Alignment Score.
